% ================================================================
%  第 12 章 第一部分的注 (原书 p.243-254, 边码 p.155-p.171)
%  源文: _work/src1/task-ch12b.md (12.26–12.50, 注 23″–44)
% ================================================================
% ---------------- p.243 ----------------
\section{怕玩}

\notehead{23$''$}{}
这个学生曾和我一起做了一整年的DEA“实习工作”,而在他与我的工作关系里,直到最后都保持着“紧绷”状态。那是一种坦率友好的关系,贯穿着一种无可怀疑的相互好感。然而可就是有这份“怯场”;这份恐惧——它真正的起因肯定不是对我这个人的畏惧,尽管看上去像是如此。要不是这个学生自己跟我谈起,我也许根本不会察觉这件事;他谈,无疑是为了或多或少“解释”一下他在这一年里工作几乎完全卡住的原因。

正如在其他学生身上发生过的那样——他们和他一样,起初对某种几何实质抓得很牢——卡住恰恰出现在要“做零件活儿”的时候:也就是把成形的陈述白纸黑字写下来,或者仅仅是把握我提供的那些陈述的意义与含义;我提议把它们当作一种语言的基石、当作“游戏规则”接受下来。“学校式的”反应几乎总是把面对“做研究”局面的学生推向:把导师传下来的那些心照不宣的“游戏规则”当作既定的、又必须服从的东西接受下来,而这些规则恰恰不是要试着说明白的,更不是要弄懂的。这些心照不宣的规则所采取的具体形式,就是语义学或计算上的“配方”,不妨说,像预科班的题解书那样(或者任何一本普通教科书)。此外,学生还期待导师给出“证明……”这种形式的任务——那是他以往经验里遇见的唯一一种数学“思考”形式。(而且我不认为大多数职业数学家、以及其他科学家的禀性有本质不同——只不过“导师”换成了“共识”:共识规定当下的游戏规则,并把它当作不可更动的既定之物。这共识还规定哪些“问题”是要去解决的;在这些问题之间,每个人都觉得自己有随意挑选的自由,甚至允许自己在工作过程中修改它们,乃至发明另一些……)。\origpage{155}我注意到,我面对有待探究的数学实质时那种完全不同的态度——因而也是面对学生时那种完全不同的态度——几乎必定会引发一种失措,而焦虑就是它的征兆之一。像一切焦虑一样,它倾向于给自己找一张面孔,投射到某个外部的“理由”上,不管那理由说得通说不通。焦虑最常见的一张面孔,恰恰就是恐惧。

这类困难在我教学生涯的第一个时期几乎没有出现,也许只有两例例外:那两次“师生”关系没超过几个星期;也许还有(我说不准)“忧伤的学生”那一例——他也许感到自己被“铆”在一个完全不启发他的题目上,而他本来完全有更换的自由。在那个(我也谈到过的)长期被某种怯场困扰的学生身上,原因显然在别处。他的工作丝毫没有卡住,恰恰相反,他对选定的主题极为自在,并在其上做出了一项有规模的基础性工作。那个时期我的学生大多是巴黎高师的旧生,他们与亨利·嘉当(Henri Cartan)的接触,已经让他们见识了另一种“接近数学”的方式。与之(可以说)相反的另一极,是我在蒙彼利埃大学任教的第二个时期:正是在一年级学生那里,我所说的焦虑最少干扰思考性的工作。在这些学生中的许多人身上,面对一种不同进路时的惊讶既不引起焦虑,也不引起封闭,相反带来的是敞开和兴致——终于可以做点有趣的事了!据我的观察,几年大学对一个学生创造性禀性的作用是彻底而毁灭性的。奇怪的是,在这方面,长年中学的作用似乎相对无害。原因也许是:大学那几年正落在这样一个年纪——我们身上与生俱来的创造力,最迟必须通过个人的工作表达出来,否则就要永远沉船,至少就智力上的创造性工作而言是如此。肯定是出于一种健康的本能,我在做学生的那些年(也是在蒙彼利埃大学),几乎不踏进课堂,把几乎全部精力都投给了个人的数学思考\origpage{156}。

% ---------------- p.244 ----------------
\section{两兄弟}

\notehead{23$'''$}{}
这个学生身上的对抗,一上来就采取了“阶级对抗”的形式:我是那个对他的数学前途拥有“生杀之权”的“老板”,可以凭自己的好恶决定……当然,事情本身只能证实这种看法,因为我很快就终止了对这个学生(已变得令人难受的)责任。这把他置于一个微妙的处境——在眼下的世道里,找一个“老板”可不是那么显而易见的事,尤其是题目已经选定的时候。在另一个学生身上,合法期待落了空,对抗采取了类似的形式。我被感受为专横的“学阀”,不能容忍那些被他视为下属的人(学生或级别较低的同事)有任何反驳。

这样一种“阶级态度”,在我与第一个时期学生的关系里,从未有过一丝一毫的表现。显而易见的原因是:在1970年以前的形势下,学生通过论文之后毫无疑问会得到一个副教授(maître de conférences)的职位,因而享有与我相同的社会地位,即“大学教授”。数字最能说明问题:1970年以前开始跟我做工作的十一个学生,一完成工作就得到了副教授职位;而在我的指导下或多或少工作过的那二十来个学生里,没有一个获得这样的职位。诚然,他们之中只有两个人有足够的动力去做国家博士论文(而且两人的论文都很出色)。

因此,在这第二个时期,某些暧昧(其深层根源始终被遮蔽着)采取了阶级对抗的形式,采取了对“老板”的不信任(这种不信任被呈现、也被感受为“本能的”“发自肺腑的”)——这就不足为奇了。在那些或多或少当过学生样子的人里,有一个与我的友好关系持续了十来年,没有出现过表面上的对抗插曲,却仍带着同一种暧昧,表现为一种不信任的态度,被“储备”在明显的善意之后。说实话,我从未被这种例行公事式的“不信任”骗过;它在我看来,尤其是这位朋友认为应当给自己找的一个理由,好让自己不去冒险越出他选定的、界限分明的领地——无论在职业生活里,还是在整个生活里都是如此。可这件事他本来可以自由地去做,没有人(顶多只有他自己!)会要他交代……

这三个例子,也是我整个教书经验中仅有的、学生(或多少像学生的人)与我之间的关系中某种暧昧以“阶级态度”表现出来的情形。当这种态度出现在大学“团体”内部的同事之间时,显得格外暧昧:在这样一个团体里,双方都享有与常人处境相比极为悬殊的特权,而这些特权使级别(以及薪金)的差别显得相对微不足道。我还注意到,一旦当事人自己也被提升到前一天他还在指责别人的那种处境,这些态度就像中了魔法一样消失了(而且理所当然!)。

此外,在我于数学界内部(常常也在外部)目睹的多数——如果不是全部——冲突局面里,我都辨认出一种类似的暧昧\origpage{157}。那些“有位置的人”,无论其级别是否合乎他们的期待(合理与否),都享有相当闻所未闻的特权,那是任何别的职业或生涯都无法提供的。那些没有位置的人,则向往同样的安全感和同样的特权(这并不必然妨碍他们对数学本身感兴趣,有时还做出漂亮的工作)。在眼下的世道里,抢位置竞争激烈,没有位置的人常被当作流浪汉对待:我不止一次感到那种默契——在乐于羞辱人的人与被羞辱的人之间;后者咽下去,缩起来。他苦涩与敌意的真正对象,并不是那个行使了权力的人,而恰恰是他自己:他自己缩了起来,并且把对方奉为那种他乐意行使的权力的持有者。乐于羞辱人的人,也是那个报复并补偿(却永远抹不掉……)一场长期承受、早已深埋并被遗忘的羞辱的人。而默许自己被羞辱的人,是他的兄弟和效仿者:他暗中羡慕对方,并在苦涩里埋下那场羞辱,以及羞辱带给他的、关于他自己的那条卑微的消息。

% ---------------- p.245 ----------------
\section{一次教学的失败(1)}

\notehead{23$^{iv}$}{}
\origpage{158}这些字句写下之后,我有机会与1970年之后的两位旧学生谈过,想同他们一起探究我在蒙彼利埃大学研究层次上教学失败的原因。他们对我说:我那种倾向——低估同化那些对我熟悉、对他们却不熟悉的技术对他们可能意味着的困难——在他们身上产生了令人泄气的效果,因为他们感到自己始终达不到我对他们的期望。此外(这件事在我看来意义还要大),他们有时感到受了挫伤:当他们已经十分接近的时候,我不是把发现的乐趣留给他们、让他们靠自己的办法去发现,而是把袖子里藏着的一个成形陈述“卖给他们”、向他们泄底。那之后,他们剩下的就只有做“练习”了(而那并不特别让他们着迷)——也就是证明那个陈述。我先前在一则注(注 21)里指出过的自己身上那份“缺乏慷慨”,正落在这一处,当时没有就此多说。主要正是这类扫兴之事,构成了我在他们两人身上——在那么出色的开端之后——研究兴趣消失这件事上的个人贡献。

我意识到,我在1970年以前并不比之后更慷慨。如果说那时我没有遇到同样的困难,那无疑是因为:那个时期来找我的那类学生,有足够的动力,甚至能从一场“长长的练习”里品出魅力——那是学习本行、并一路学到一大堆东西的机会;同样,对于一个我“泄了底”的起步陈述,他们也能靠自己的办法从中抽出一大群远远超出第一个的陈述。当我换了教学地点,我在向新学生提议的思考主题上做了必要的调整,办法是选择那些可以被直接直觉把握的数学对象,不依赖任何技术行囊。但这必要的调整本身并不够,因为禀性上的差别(我的新学生与从前的学生相比)比单纯的行囊差别还要重要。这还印证了前面(第25段开头)说过的一点:我在“导师”这个角色上有某种不足;这不足在我教学生涯的第二个时期,比在第一个时期显露得强烈得多。

\section{$\varnothing$}

% ---------------- p.246 ----------------
\notehead{23$^{v}$}{}
\origpage{159}这种差别有一个格外醒目的征兆,出现在“外国人事件”之际——我曾有机会谈到过(第24节)。当时,许多与我完全陌生的人向我表达了同情,而我记不起1970年之前的哪个学生想到过以这种方式表示什么,更不用说在我投身的那场行动中向我提供任何帮助。相反,在我看来,第二个时期的学生或旧学生里,没有一个不曾向我表达同情与声援,还有几个积极参加了我在当地开展的运动。在这个狭小圈子之外,1945年法令那件事也在许多仅仅知道我名字的大学生中间激起了一定的波澜;在我被传唤的那一天,有不少人来到法院,以示声援。这后一个情形还提示:我所察觉到的“之前”与“1970年之后”学生态度上的差别,也许较少表达他们与我之间关系的差别,而更多表达的是一种心态的差别。显然,我“之前”的那些学生已经成了重要人物,而要让重要人物动心,是需要很多条件的……但我1970年离开IHÉS、投身战斗行动那一幕,似乎表明并不只是这一点。那是个他们中间还没有谁太像重要人物的时刻,然而我记不起他们中有谁对我所投身的活动表现出丝毫兴趣。我倒认为,那一定让他们所有人,无一例外地感到不自在。这仍然合乎“心态差别”的方向,但这不能只归因于社会地位的差别。

\section{$\varnothing$}

\notehead{24}{}
我要谈的这套伦理,同样适用于其他一切围绕研究活动形成的圈子:在那里面,能不能让人知道自己的结果、并从中得到承认,对每个成员的社会地位都是“生死”问题,甚至是他作为该圈子成员的“存亡”问题,而这对本人及其家庭意味着什么,也就可想而知。

\section{职业道德的共识——以及信息的控制}

\notehead{25}{}
除了与迪厄多内的那次谈话,在我作为数学家的生涯里,我想不起还有哪一次我参与或旁听的谈话\origpage{160},谈到了这一行的伦理、谈到同行之间关系里的“游戏规则”。(这里不算七十年代初围绕“生存与生活”运动发生的那些关于科学家与军事机构合作的讨论。那些讨论并不真正涉及数学家彼此之间的关系。再说,我在“生存与生活”里的许多朋友,包括谢瓦莱和盖吉,都觉得我当时——尤其在初期——如此强调这个我格外敏感的问题,使我偏离了更紧要的日常现实,恰恰就是眼下这份反思所考察的那类现实。)在我与学生之间,从未谈到过这些事。我相信,那份默契只限于一条规则:不把别人的想法——只要自己有机会知道——说成自己的。我觉得这条共识自古就有,直到今天在任何科学圈子里都没有被质疑过。可是,若缺少另一条补充规则——保证每个研究者都有可能让人知道自己的想法与结果——第一条规则就只是一纸空文。在今天科学界,处在声望与权力位置上的人,把持着对科学信息的任意支配权。
% ---------------- p.247 ----------------
在我从前熟悉的那个圈子里,这种把持已不再被迪厄多内所说的那种共识所节制——而那种共识,也许从未存在于他充当代言人的那个小圈子之外。处在权力位置上的科学家,实际上能收到他认为值得收到的一切信息(往往还超出);而对于其中很大一部分,他有权阻止其发表,同时自己留着这份信息的好处,把它打发为“没意思”“多少已经知道了”“平凡”之类……这种局面我在注(27)里还要再谈。

\section{$\varnothing$}

\notehead{26}{}
布尔巴基的“创始成员”是亨利·嘉当、克劳德·谢瓦莱、让·德尔萨特(Jean Delsarte)、让·迪厄多内、安德烈·韦伊。他们全都健在,只有德尔萨特例外——五十年代,他未及天年便被带走了;那个年头,这一行的伦理还普遍受到尊重。

重读本文时,我曾想删掉这一段:在这里,我可能给人一种印象,仿佛在颁发“正直”(或不正直)证书,而当事人对此毫不在乎,这事也不该由我来做。这一段可能引起的保留,肯定是有道理的。\origpage{161}可我仍旧留着它,为的是证词的真实,也因为这一段确实如实复现了我的感受,哪怕这些感受并不合宜。

\section{“年轻人的势利”,或纯粹性的捍卫者}

\notehead{27}{}
罗尼·布朗(Ronnie Brown)向我转述了J.H.C.怀特黑德的一句话(他是怀特黑德的学生),说的是“年轻人的势利:他们以为一条定理平凡,只因为它的证明平凡”。我从前的许多朋友,不妨把这句话好好想一想。这种“势利”如今绝不限于年轻人;我认识不止一位有声望的数学家,常这么做。我对此格外敏感,因为我在数学上(别处也一样……)做得最好的那些东西——我引入的、在我看来最丰饶的概念与结构,以及我靠耐心而执拗的工作从中拔出的那些本质性质——全都落进了“平凡”这个评语里。(如今这些东西,若作者还不是个名人,哪一件都很难被《科学院通报》(CR)当作一篇短讯接受!)我这一辈子作为数学家的抱负——不如说是我的激情与快乐——始终是找出那些显然的东西;在这本书里(包括眼前这一章引论性的文字里……),这也同样是我唯一的抱负。决定性的东西,常常首先就是看见那个尚未被看见的问题(不管答案是什么,也不管答案是否已经找到),或者提出一条陈述(哪怕是猜想性的),把一种尚未被看见、或尚未被理解的局面概括并容纳进来;这条陈述若被证明了,那么证明平凡与否无关紧要——那完全是附带的事——甚至一条仓促的、暂用的证明后来被证明是错的,也无关紧要。怀特黑德所说的势利,是那种厌倦了的享乐者的势利:他要先确认一瓶酒贵得惊人,才肯去欣赏它。这些年里,不止一次,旧日的激情一阵阵回来抓住我,我把自己最好的东西捧出去,却看它被那份自得打了回来。我感到一种至今仍活的痛;一份快乐落了空——可我并没有因此流落街头,而且幸而我并不需要替自己写的一篇文章找地方发表。

% ---------------- p.248 ----------------
\origpage{162}怀特黑德所说的势利,当它由一个有权者施于一个任其摆布的研究者身上时,就不只是一种对事物之美的麻木或封闭,而是一种滥用权力、一种不诚实:那位研究者完全在他的掌握之中,他尽可随意吸收、使用其想法,同时以“显然”“平凡”因而“没意思”为借口,堵住它们的发表。我这里甚至没有想到通常意义上的剽窃那种极端情形——在数学界,那想必还极罕见。可是从实际角度看,对为此付出代价的研究者来说,局面是一样的;而使这局面成为可能的内心态度,在我看来也差不太多。它只是更舒服些:因为它伴着一种对他人无限优越的感觉,伴着良心安然与内心满足——那是一个自命为数学之不可触犯的纯粹性的不妥协捍卫者的人。

\section{$\varnothing$}

\notehead{28}{}
写前面那几页时,我起初在“把心里的东西全倒出来”的欲望与一份保留、谨慎的顾虑之间分裂。于是我就停在差不多、大约莫的地方——这肯定是我那份不适、“我什么也没学到”之感的主要原因。自从写下那几行承认不适的文字之后,我把这几页重写了两遍——它们让我心里不满意;我把自己更清楚地放了进去,也把事往更深处推。一路下来,我确实终于“学到了点东西”;我也相信,与此同时我摸到了某件要紧的东西,它超出了这个具体个例,也超出了我这个人。

\section{$\varnothing$}

\notehead{29}{}
我这里说的,是在数学里、或在另一项纯智力活动里那种强烈而长久的投入。反过来,这样一份激情的展开——它可以是重新认识我们身上一种被遗忘的力量的方式,也是与一种不肯就范的实质较量的机会,并且一路下来,用某种真正属于我们自己的东西去更新、丰富我们的身份感——这样一种展开,很可能是内心路程上的一个重要阶段,是成熟过程中的一个重要阶段。

\section{$\varnothing$}

\notehead{30}{}
这几年,接力的是我的孩子们:由他们去教一个有时不大情愿的学生,人生的种种奥秘……

\section{$\varnothing$}

\notehead{31}{}
\origpage{163}我这里想到的是求知欲的“阳”的形式——那探问、发现、为显现者命名的形式……正是由于被命名,出现的知识才变得不可逆转、不可抹去(哪怕它后来被埋掉、被遗忘,不再活跃……)。求知欲的“阴”的、“女性”的形式,在于一种敞开、一种接纳,在于无声地迎候一种出现在我们存在更深层——思想到不了的地方——的知识。
% ---------------- p.249 ----------------
这样一种敞开、一种突然到来、一时间抹去一切冲突痕迹的知识的出现,仍像一份恩典:它触到深处,而它可见的效果也许转瞬即逝。不过我猜,这种在我们一生中某些稀有时刻如此到来的、无言语的知识,同样不可抹去;它的作用甚至延续到我们对它的记忆之外。

\section{炉中百铁,或:干熬无用!}

\notehead{32}{}
我还在做泛函分析的时候,也就是直到1954年,我有时会没完没了地死咬一个问题,解不出来;那时我已经没有新想法,只是绕着旧想法的圈子打转,而那些想法显然已经“咬不住”了。整整一年都是这样,尤其是拓扑向量空间里的那个“逼近问题”;它要再过二十年才被解决,用的完全是另一路的方法,以我当时所处的境地,那些方法只能与我擦肩而过。当时驱动我的不是欲望,而是一种固执,以及对自己内心所发生之事的不了解。那是难受的一年——我一生中唯一一段做数学变得难受的时期!我需要这段经历才明白:“干熬”没有用——一件工作一旦到了停顿点,而且只要停顿被察觉,就该转到别的事上去;至于悬置的那个问题,等到更合宜的时候再回来。这个时刻几乎总是很快就出现——问题自己成熟起来,我并没有做出碰它的样子;靠的只是兴致勃勃地做那些看来与它毫无关系的问题的功效。我确信,当时若我硬撑下去,哪怕十年也一事无成!从1954年起,我在数学上养成了一个习惯:炉子里总同时烧着许多块铁。我一次只打一块,可由于一种不断更新的奇迹\origpage{164},我在一块上做的工作也惠及所有其余的,它们各自等着自己的时辰。与我第一次接触冥想时也一样,而且完全不是我刻意为之:随着反思继续,有待审视的灼热问题,数目一天天增多。

\section{$\varnothing$}

\notehead{33}{}
这并不意味着:工作中纸(或黑板,那不过是它的一种变体!)不在场的那些时刻,在数学工作里就不重要。尤其如此的是那些“敏感时刻”:一种新的直觉刚刚出现,这时要做的是以一种更整体、更直觉的方式去“认识”它,而不是做“零件活儿”——而反思的这个不定形阶段,正是为那件事作准备。在我身上,这类思考主要在床上或散步时进行;我觉得它只占用于工作的全部时间中一个相对不大的部分。同样的观察也适用于我迄今所从事的冥想工作。

\section{无力的拥抱}

\notehead{34}{}
“拥抱”这个词在我这里绝不是单纯的比喻;日常语言在此映出的是一种深层的同一。人们不无道理地会说:这么说,“没有惊叹的拥抱无能为力”就不对了——若按字面如此,大地就该人口凋零,或者干脆荒无人烟。极端的例子是强奸:惊叹当然不在场,可有时被强奸的女人身上仍孕育出一个生命。由这样的拥抱生下的孩子,肯定免不了带着那个印记;那印记会是他分得的那一份“包裹”的一部分,该由他自己去承担。可这并不妨碍:一个新生命确实被孕育、被生了下来;确实有过创造,那是力量的标志。也确实如此:我见过某个满身自得的数学家找到并证明了漂亮的定理,那正是拥抱不乏力气的标志!但同时也不假:若某位数学家的生活被他的自得闷死(在某种程度上,我自己生活中某个时期也曾如此),那么他与数学的这些拥抱所结的果子,对他、对别人都不是什么福分。从强奸中生下的孩子的父亲与母亲,也可以这样说。我说“无力的拥抱”,首先指的是这样一件事:在那个自以为在创造的人身上,无力生出更新——他只造出一个产品、一件他之外的东西,在他自身里没有深的回响;
% ---------------- p.250 ----------------
\origpage{165}这产品不但不解放他、不在他里面造出和谐,反而把他更紧地捆在他身内那份他作囚徒的自得上,那自得不断催他去生产、再生产。这是深层上的一种无力,藏在“创造力”的表象背后——而那归根到底不过是无缰的生产力。

我也有充分机会看出:自得、无惊叹之能,其性质是一种真正的盲,是自然感受力与嗅觉的堵塞;这堵塞即使不是全面而永久的,至少在某些具体情形里是明显的。在这种状态里,某位有声望的数学家有时会在恰恰他擅长的那些事情上,显得比最死心眼的小学生还蠢!在另一些场合,他会做出技术精湛的奇迹。可我怀疑,他是否还有能力去发现那些单纯而显然、有力量更新一门学科或一门科学的东西。那些东西离他太低、太远,他已经不屑于去看!要看别人不屑于看的东西,需要一份他早已失去、或被他放逐的天真……近二十年间数学产出惊人地增长,新结果的纷繁令人无所适从,连只想稍微“跟上”的数学家都被淹没;而与此同时(就我从四处传来的回声所能判断的而言),我所多少熟悉的那些大思考主题,却几乎没有真正的更新、没有大幅度的转化(不只是积累式的转化)——这肯定不是偶然。更新不是量的事情;它与投入的量无关——那种量可以用若干“数学家人日”来计量,即某些“水平”的某些数学家花在某个题目上的日数。一百万个数学家人日,也无力生出像零那样孩子气的东西——而零更新了我们对数的感知。只有天真拥有这份力量;它的一个可见标志,就是惊叹……

\section{$\varnothing$}

\notehead{35}{}
这份“天赋”不是谁的特权,我们全都生来带着它。当它在我身上显得不在时,那是我自己把它赶走的;要不要重新接纳它,只在于我。在我身上,或在某某身上,这份“天赋”的表达方式与另一个人不同,不那么有感染力,也许不那么不可抗拒,但它照样在场;我也说不出它是否不那么有力。

\section{$\varnothing$}

\notehead{36}{}
\origpage{166}这种对美的细腻感受力,在我看来与一样东西紧密相连;
% ---------------- p.251 ----------------
我曾以“严要求”(对自己的)或“严格”(取其完整意义)之名谈到它,把它描述为“对我们自身中某种细腻之物的注意”,对所要探问之物的理解品质的注意。对一件数学东西的理解品质,离不开对该物特有之“美”的一种或深或浅、或全或欠的感知。

\section{$\varnothing$}

\notehead{37}{}
我想,几乎不必补充:这项长久的工作日复一日地显出许多别的东西,远不止我刚才以碑铭般的简省形式交出的那个“结果”。冥想的工作,与由某个打算考察的具体问题所驱动的数学工作,并无不同。所走那条路(它通向、或不通向对最初问题或多或少的澄清)上的种种曲折,往往比最初的问题或“最终结果”更有意思。

\section{$\varnothing$}

\notehead{38}{}
这些注其实是写给……的那封长信的继续,那封信成了本书第一章。它们用打字机打出来,是为了让这位旧日的朋友能读,也为了另外两三位(尤其是罗尼·布朗),我以为他们可能会感兴趣。再说,那封信从未得到回音,收信人也没有读它:将近一年之后(我问过他是否确实收到),他由衷地惊讶——我怎么会有一刻以为他会读它,鉴于人们该期待我做的那类数学……

\section{$\varnothing$}

\notehead{39}{}
这一时期,除别的之外,也正是《穿越伽罗瓦理论的长征》的时期;在《纲领概要》里(第3段:“与童画相联系的数域”)谈到过它。

\section{来访}

\notehead{40}{}
对这个梦所做的工作,是一封英文长信的内容;收信人是一位朋友兼同行,他前一天在我家匆匆过了一趟。做梦者用来从一片表面的虚无中唤出这个逼真得惊人的梦的那些材料,有些显然取自那位亲爱的朋友来访的短暂片段——我已有近十年没再见过他。\origpage{167}于是,在工作头一天,而且与我过去的经验相违,我竟以为可以由此断定:我做的这个梦关系到我的这位朋友,多过关系到我——该做这个梦的是他,不是我!这是一种回避梦境消息的办法;而那消息(凭我过去的经验,我本该一开始就知道)只关系到我,不关别人。在这第一阶段——浮浅的——工作之后的夜里,我终于意识到了这一点;第二天,我在同一封信里接着做了下去。自那封值得记住的信之后,我再没有收到这位朋友的任何音讯——他是我有过的最亲近的朋友之一。

% ---------------- p.252 ----------------
这项工作,是唯一一次以信件形式(而且用的还是英文)成形的冥想,因此我手里不再有它的书面痕迹。这段插曲让我格外震动——在众多别的例子当中,它们都表明:任何越出某层门面的工作迹象,任何把极简单的事实——而人们通常认为有义务对之视而不见——带到光天化日之下的工作,是多么在别人身上引起不适与恐惧。这一点我在后面还要谈(见第47段《孤独的历险》)。

\section{克里希那穆提,或变成羁绊的解放}

\notehead{41}{}
若说我从这次阅读中得到的只有一套词汇,以及一种把它据为己有、最终又顺理成章地用它取代现实的倾向,那是不准确的。我拿到手的第一本克里希那穆提(Krishnamurti)的书之所以那么震动我(而且我还只来得及读其中几章),是因为他说的东西彻底掀翻了我许多自以为当然的事;我立刻看出,那些都是老生常谈,一直是我所呼吸的空气的一部分。同时,这次阅读第一次把我的注意力引向一些分量很重的事实,尤其是对现实的逃避——它是心智最强大、最普遍的制约之一。这给了我一把关键的钥匙,去理解那些此前一直不可理解、并因此(在我五六年之后发现冥想之前,我并未察觉这一点)滋生焦虑的处境。我立刻就能在我周围各处看到这种逃避的现实。这解开了某些焦虑,却没有改变任何本质的东西,\origpage{168}因为我只在他人身上看到那份现实,同时(理所当然地)以为它在我身上并不存在,以为我总之是那条证实规则的例外(而且对这么一个确实了不起的例外,我没有再问任何别的问题)。事实上,我对别人、对自己都毫无好奇。这把“钥匙”只有在被深入之欲驱动的人手里才能开锁。在我手里,它成了驱邪和做姿态。

直到1974年初,我才第一次认清:我生活里的那种毁灭一步步跟着我,不可能只来自别人;我身上有某种东西在吸引它、喂养它、让它延续。那是一个谦卑与敞开的时刻,适宜于一次更新。可那次更新当时仍是边缘的、短暂的,因为缺少深入的工作。那个“我身上的某种东西”还很模糊。我看得出那是爱的匮乏;可是,要有一项工作去更贴近地界定:爱的匮乏在我身上何处、如何发生,它如何显现,它具体造成了什么后果,等等——这样的念头,既不可能来自我至今认识的任何圈子或任何人,也不可能来自克里希那穆提。(恰恰相反,K.喜欢强调一切工作的徒劳,他自动把工作等同于“我”的“成为之饥”。)就这样,我拿一份借来的“智慧”当唯一的罗盘,看不见还有别的事可做,只好耐心等着“爱”像圣灵的恩典一样降到我身上。

然而,我刚在一个浪的深处学到的那点卑微真相,已经激起了一股强劲的新能量之浪,可与两年半之后承载我第一次投入冥想的那股相比。这份能量当时并没有完全闲置。几个月后,我因一场凑巧的事故动弹不得,它承载了一份(写下来的)反思:我一生中第一次审视那套世界观——它一直是我与他人关系未曾言说的基础,来自我的父母,尤其来自我母亲。当时我非常清楚地看出:这套看法已经破产,它无法说明人与人之间关系的实际,也无法促成我这个人以及我与他人关系的舒展。这份反思仍带着
% ---------------- p.253 ----------------
“克里希那穆提风格”的印记,也带着克里希那穆提式的禁忌——禁止一切朝向理解的真正工作。可它毕竟把几个月前生出的一份认识变得可触而不可逆转——\origpage{169}那认识起初一直模糊、游移。那份认识,当时世上任何一本书、任何别的人都不可能带给我。

要够得上冥想的品质,这份反思最缺的,是对我自己这个人、对我关于自己的看法的注视——而不只是对世界观的注视,对一个公理系统的注视,而在那系统里我并不真的“有血有肉”地出场。它还缺另一样:在当下、在反思进行的那一刻对我自己的注视(而那份反思还停留在真正工作的此岸);那样一种注视,本来会让我既察觉一种借来的风格,也察觉我对这些笔记的文学面貌有某种自得,也就是说,缺少自发与真实。尽管它如此不足,尽管它对我与他人关系的直接影响相对有限,这份反思在我看来仍是一个阶段——考虑到起点,这大概是必需的——通向两年后该发生的那次更深的更新。正是在那时,我终于发现了冥想——同时发现这第一个未曾料到的事实:关于我自己这个人,是有东西可发现的;而那些东西几乎完全决定了我一生的走向,以及我与他人关系的性质。

\section{有益的剥离}

\notehead{42}{}
那个“有冲击力”的事件,是1969年底的发现:我觉得自己属于的那个机构,部分经费来自国防部的款项——这与我基本的公理不相容(今天也仍然不相容)。这个事件是一整串别的(一个比一个更有揭示力!)事件中的第一个;它们的结果是我离开高等科学研究所(IHÉS),并由此一环扣一环地彻底换了圈子与投入。

在IHÉS的英雄年代,迪厄多内和我是它仅有的成员,也是仅有的为它在科学界赢得信誉与听众的人:迪厄多内靠主编《数学出版物》(头一卷早在1959年就出了,即莱昂·莫尚(Léon Motchane)创办IHÉS的次年),我靠《代数几何讨论班》。最初那些年,IHÉS的存在极为飘摇:经费没有着落(靠几家\origpage{170}充当赞助人的公司的慷慨),唯一的场所是巴黎蒂耶尔基金会(Fondation Thiers)借出的一间厅(它显然很不情愿),专供我讨论班开讲的日子\fnote{5}{}{IHÉS最近为纪念其成立二十五周年出的一本小册子(承尼科·凯珀(Nico Kuiper)好意寄我一本),对这样艰难的开端只字不提——也许认为它们配不上这个场合的庄重吧;去年,这个周年被大张旗鼓地庆祝了一番。}。我多少觉得自己,连同迪厄多内,是我所属这个机构的“科学”共同创办人;我还真打算在那儿过完这辈子!我终于对IHÉS有了很强的认同;而我的离开(由于同事们的漠然),当时被体验为从另一个“家”被剥离,后来才显出是一次解放。

回过头看,我意识到自己身上想必早有更新的需要,我说不清是从什么时候起。我离开IHÉS的前一年,我的精力投入突然翻了个面:丢下那些昨天还烫手的工作、那些最让我着迷的问题,转而去投身生物学(是在一位生物学家朋友米尔恰·杜米特雷斯库(Mircea Dumitrescu)的影响下)——这肯定不是单纯的巧合。我投身其中,是打算在IHÉS内部做
% ---------------- p.254 ----------------
\origpage{171}长期投入的(这也符合该机构多学科的宗旨)。这肯定只是给一种深得多的更新需要找了个出口;那种更新在IHÉS那种“科学蒸笼”的气氛里无从完成,而是在我已提及的那场“觉醒的瀑布”中发生的。共有七次,最后一次在1982年。“军费”那段插曲是天赐的,因为它触发了这些“觉醒”中的第一次。国防部,以及我在IHÉS的旧同事们,最终都配得上我全部的感激!

\section{$\varnothing$}

\notehead{43}{}
“我写的那部诗作”里包含许多我第一手知道的东西;它们今天在我生活中、在一般的“生活”中,显得与它写成之时同样重要——那时我是打算把它发表的。我之所以没有发表,主要是因为后来意识到:那形式被一种刻意“作诗”的意图所累,以至于它整体构思过于搭建,许多段落缺少自发,有些地方甚至僵硬或浮肿得难受。这种有时浮夸的形式,是我当时禀性的映照:在那里面,常常确实是“老板”在领舞——笨重得很,那是不用说的……

\section{$\varnothing$}

\notehead{44}{}
不用说,我在这里是把这样一个假设撇开的——至少可以说,它并非不可能:一场原子战争,或另一场同类的“乐事”,突然闯来,把那个叫做“数学”的集体游戏一举粗暴地终结,顺带还终结别的好些东西……
